\chapter{Introdução}
\label{intro}

No âmbito da unidade curricular de Administração de Sistemas, foi proposto o desenvolvimento de um projeto orientado à automatização da configuração de serviços num servidor Linux. O enunciado define um conjunto de funcionalidades técnicas centradas na gestão de DNS, integração com serviços Web, administração de partilhas de rede, mecanismos de backup e reforço da segurança do sistema, privilegiando sempre a execução automática por scripts.

Este trabalho surge da necessidade de reduzir tarefas manuais repetitivas que, em contextos de administração de sistemas, tendem a aumentar a probabilidade de erro, dificultar a normalização de configurações e tornar a manutenção mais morosa. Neste sentido, a automação é assumida como estratégia principal para melhorar a eficiência operacional, garantir maior consistência entre configurações e facilitar a reprodução dos procedimentos técnicos em diferentes momentos de execução.

O objetivo geral consiste em implementar uma solução capaz de configurar e gerir, de forma automatizada, um conjunto de serviços de administração de sistemas em Linux, com validação em ambiente virtualizado. O desenvolvimento é realizado com recurso a, pelo menos, uma máquina servidor Linux e uma máquina cliente para testes, ambas em VirtualBox, em conformidade com os requisitos do projeto.

Ao longo do relatório serão apresentados, de forma progressiva, o enquadramento técnico das opções adotadas, as etapas de implementação e os resultados de validação. A estrutura do documento segue uma lógica de desenvolvimento que inclui a descrição detalhada dos procedimentos técnicos, a análise crítica das soluções implementadas e a reflexão sobre os desafios enfrentados durante o processo de automação.
